% Generated draft: requires mathematical and source review.
\documentclass[11pt]{article}
\usepackage[margin=25mm]{geometry}
\usepackage{fontspec}
\setmainfont{Latin Modern Roman}
\usepackage{amsmath,amssymb,mathtools,unicode-math}
\setmathfont{Latin Modern Math}
\usepackage{enumitem,needspace,xurl,hyperref}
\setlength{\parindent}{0pt}
\setlength{\parskip}{5pt}
\setlength{\emergencystretch}{3em}
\newcommand{\N}{\mathbb{N}}
\newcommand{\Z}{\mathbb{Z}}
\newcommand{\Q}{\mathbb{Q}}
\newcommand{\R}{\mathbb{R}}
\newcommand{\C}{\mathbb{C}}
\newcommand{\F}{\mathbb{F}}
\newcommand{\A}{\mathbb{A}}
\newcommand{\PP}{\mathbb{P}}
\newcommand{\E}{\mathbb{E}}
\newcommand{\eps}{\varepsilon}
\newcommand{\dd}{\,\mathrm{d}}
\begin{document}
\section*{UnsolvedMath ID: 20002289}
\textbf{Source problem code:} catalogue entry
\subsection{Definitions and mathematical statement}
Koenigsmann's Question 5 asks whether every inclusion of fields
\[
 F\subseteq E,\qquad F\equiv\Q,\quad E\equiv\Q
\]
in the language of rings satisfies
\[
 \{b\in E:b\text{ is algebraic over }F\}=F.
\]
The inclusion is not assumed elementary. Equivalently \(F\) should be relatively algebraically closed in \(E\); the source's notation \(F^{\mathrm{alg}}\cap E=F\) uses a common algebraic overfield. Its reported preliminary fact is relative quadratic closedness, meaning no new quadratic roots in \(E\), not that \(F\) itself is a quadratically closed field.

The aid reports a positive full answer from Dittmann's Theorem 5.2 and Corollary 5.3. For each fixed degree \(n\ge1\), let
\[
 R_n(a_0,\ldots,a_{n-1})\quad\Longleftrightarrow\quad
 \forall x\ \left(x^n+\sum_{i=0}^{n-1}a_ix^i\ne0\right).
\]
The named theorem existentially defines monic rootless coefficient tuples over global fields, in particular over \(\Q\), and the corollary gives the asserted relative algebraic closedness for these models. The summary does not state uniformity of one formula across all fields or all degrees; no such uniformity is added.

The hypothesis bridge checked by the aid is that the canonical prime-field copy of \(\Q\) is elementary in every model of \(\operatorname{Th}(\Q)\). Rational parameters are definable from ring constants, so their formulas transfer with that canonical identification. Consequently the fixed-degree existential rootlessness definitions transfer to \(F\) and \(E\), and existential witnesses persist from \(F\) to \(E\). This explains the use of the known result without replacing the given inclusion by an elementary embedding.

The aid separately proposes a bounded-degree organizational criterion. For a complete theory \(T\) of fields and integer \(d\ge1\), the following are reported equivalent:
\begin{enumerate}
\item Every inclusion \(F\subseteq E\) of models of \(T\) has no \(b\in E\setminus F\) algebraic over \(F\) of degree at most \(d\).
\item For every \(1\le n\le d\), every monic degree-\(n\) polynomial with coefficients in \(F\) that has no root in \(F\) still has no root in \(E\), for every such inclusion.
\item For each \(1\le n\le d\), the predicate \(R_n\) is equivalent modulo \(T\) to an existential formula.
\end{enumerate}
This criterion is the low-confidence candidate contribution; it is not attributed as the source of Dittmann's known full solution.

The original adds that a negative answer for some inclusion would show \(\Z\) is not existentially definable in \(\Q\). The reported positive answer does not establish that \(\Z\) is existentially definable there; that converse is not inferred.
\subsection{Short English statement}
The question asks whether a subfield model of the rationals can acquire new algebraic elements inside another such model without an elementary inclusion. The aid reports Dittmann's positive solution through existential rootlessness predicates and checks the prime-field hypothesis bridge. Its separate bounded-degree criterion and the original implication about defining the integers are kept logically distinct.
\subsection{Sources}
\begin{itemize}
\item Local UnsolvedMath registry, permanent ID 20002289, exact source code catalogue entry; canonical \texttt{aim-logic-notes.json}, \texttt{notes[64]}, tag \texttt{question}; AIM Workshop Problem Lists, Logic, Question 5 (Koenigsmann). Created and updated 14 August 2026.
\item \emph{Definability and Decidability Problems in Number Theory}, AIM workshop 9--13 September 2013, moderated by T. Scanlon, notes by J. Demeyer, according to supplied verification; \url{https://aimath.org/pastworkshops/definabilityinntproblems.pdf}. The PDF was not accessed.
\item Dittmann, Theorem 5.2 (existential fixed-degree rootlessness over global fields) and Corollary 5.3 (the positive answer), as identified by the aid. No full title, publication date, identifier, or separate URL is supplied.
\item Attempt 1: \texttt{full\_solution}, problem status at run \texttt{solved}, rigor label \texttt{proved}, novelty confidence \texttt{low}; classification \texttt{SOLVED-IN-LITERATURE}, registry status \texttt{solved}. Candidate: bounded-degree criterion, explicitly separate from the known solution. These are dated source labels, not independent certification.
\item Ulam AI / UnsolvedMath Contributors; CC BY 4.0. Machine-generated research aid, not peer reviewed; independently verify before citation.
\end{itemize}
\end{document}
